%%%%%%%%%%%%%%%%%%%%%%%%%%%%%%%%%%%%%%%%%
% a0poster Portrait Poster
% LaTeX Template
% Version 1.0 (22/06/13)
%
% The a0poster class was created by:
% Gerlinde Kettl and Matthias Weiser (tex@kettl.de)
% 
% This template has been downloaded from:
% http://www.LaTeXTemplates.com
%
% License:
% CC BY-NC-SA 3.0 (http://creativecommons.org/licenses/by-nc-sa/3.0/)
%
%%%%%%%%%%%%%%%%%%%%%%%%%%%%%%%%%%%%%%%%%

%----------------------------------------------------------------------------------------
%	PACKAGES AND OTHER DOCUMENT CONFIGURATIONS
%----------------------------------------------------------------------------------------

\documentclass{a0poster}
\usepackage[pass,paperwidth=34in,paperheight=22in]{geometry}

\usepackage{multicol} % This is so we can have multiple columns of text side-by-side
\columnsep=100pt % This is the amount of white space between the columns in the poster
\columnseprule=3pt % This is the thickness of the black line between the columns in the poster

\usepackage[svgnames]{xcolor} % Specify colors by their 'svgnames'
\usepackage{graphicx} % Required for including images
\usepackage{booktabs} % Top and bottom rules for table
\usepackage{mathdots}
\usepackage[font=small,labelfont=bf]{caption} % Required for specifying captions to tables and figures
\usepackage{amsfonts, amsmath, amsthm, amssymb} % For math fonts, symbols and environments
\usepackage{wrapfig} % Allows wrapping text around tables and figures
\makeatletter
\renewenvironment{abstract}{%
    \if@twocolumn
      \section*{\abstractname}%
    \else %% <- here I've removed \small
      \begin{center}%
        {\bfseries \LARGE\abstractname\vspace{\z@}}%  %% <- here I've added \Large
      \end{center}%
      \quotation
    \fi}
    {\if@twocolumn\else\endquotation\fi}
\makeatother

\usepackage{sectsty}
\sectionfont{\fontsize{60}{60}\selectfont}
\subsectionfont{\fontsize{50}{60}\selectfont}
\subsubsectionfont{\fontsize{40}{48}\selectfont}

\begin{document}

%----------------------------------------------------------------------------------------
%	POSTER HEADER 
%----------------------------------------------------------------------------------------

% The header is divided into two boxes:
% The first is 75% wide and houses the title, subtitle, names, university/organization and contact information
% The second is 25% wide and houses a logo for your university/organization or a photo of you
% The widths of these boxes can be easily edited to accommodate your content as you see fit

\begin{minipage}[c]{\linewidth}
\begin{centering}
\vspace{-1cm}
\veryHuge \color{NavyBlue} \textbf{Data Driven Reduced Order Modeling} \color{Black}\\[1cm] % Title
%\Huge\textit{Subtitle}\\[2cm] % Subtitle
\Huge Aryn Harmon$^1$ and Jason Turner$^2$\\
\huge University of Illinois at Urbana-Champaign, Champaign, Illinois, U.S.A.\\[0.4cm] % University/organization
\Large $^1$ \texttt{arynh2@illinois.edu}, $^2$ \texttt{jasonet2@illinois.edu}\\
\end{centering}
\end{minipage}

\vspace{1cm} % A bit of extra whitespace between the header and poster content

%----------------------------------------------------------------------------------------

\begin{multicols}{3} % This is how many columns your poster will be broken into, a portrait poster is generally split into 2 columns

%----------------------------------------------------------------------------------------
%	ABSTRACT
%----------------------------------------------------------------------------------------

\color{Navy} % Navy color for the abstract

\hrulefill
\vspace{0.5cm}
\begin{abstract}
\LARGE

\textsc{Reduced Order Modeling} (ROM) can be applied to high-dimensional problems to reduce the computational cost of finding solutions. A high fidelity model was created for several partial differential equations of physical significance. Using data from these models, accurate reduced order models were constructed.

\end{abstract}
\hrulefill

%----------------------------------------------------------------------------------------
%	OBJECTIVES
%----------------------------------------------------------------------------------------

\color{Black} % DarkSlateGray color for the rest of the content

\section{Main Objectives}
\LARGE
\begin{enumerate}
\item Make online computation possible in computationally challenging problems.
\item Speed up analysis of problems with many degrees of freedom.
\item Produce an understanding of physical processes which are currently obscured by complexity.
\end{enumerate}

\vspace{1cm}
\begin{minipage}{.5\linewidth}
  \centering
  \includegraphics[width=.95\linewidth]{karmansquare}
  \captionof{figure}{Atmospheric von Karman Vortex Street \cite{environmental}}
\end{minipage}
\begin{minipage}{.5\linewidth}
  \centering
  \includegraphics[width=.95\linewidth]{karmansim}
  \captionof{figure}{Simulated von Karman Vortex Street \cite{ruggirello_2003}}
\end{minipage}

%----------------------------------------------------------------------------------------
%	INTRODUCTION
%----------------------------------------------------------------------------------------

\color{Black} % SaddleBrown color for the introduction
\section{Introduction}
\LARGE
\begin{itemize}
  \item Linear Convection and Diffusion Equation describes movement and diffusion of a wave in one dimension.
  $$\frac{\partial u}{\partial t} + c \frac{\partial u}{\partial x} = \nu \frac{\partial^2 u}{\partial x^2}$$
  $$\vec{u}^{\,n+1} = \mathcal{L}(\vec{u}^{\,n}) = \mathrm{\textbf{B}} \vec{u}^{\,n}$$
    
  \item Viscous Burger's Equation
  $$\frac{\partial u}{\partial t} + u \frac{\partial u}{\partial x} = \nu \frac{\partial^2 u}{\partial x^2}$$
  $$\vec{u}^{\,n+1} = \mathcal{L}(\vec{u}^{\,n}) + \mathcal{N}(\vec{u}^{\,n})$$
  
  \item Black-Scholes Model prices European options.
  $$\frac{\partial w}{\partial t} = rw - rx \frac{\partial w}{\partial x} - \frac{1}{2} v^2 x^2 \frac{\partial^2 w}{\partial x^2}$$
\end{itemize}

%----------------------------------------------------------------------------------------
%	MATERIALS AND METHODS
%----------------------------------------------------------------------------------------

\section{Methods}
\LARGE
%---------------------------------------------------
\subsection{Full Order/High Fidelity Model}
\subsubsection{Linear Convection and Diffusion Equation}
Finite Difference Scheme:
$$\frac{\partial u}{\partial x} = \frac{u_{i+1} - u_{i-1}}{2 \Delta x} + \mathcal{O}(\Delta x^2)$$
$$\frac{\partial^2 u}{\partial x^2} = \frac{u_{i+1} - 2u_i + u_{i-1}}{\Delta x^2} + \mathcal{O}(\Delta x^2)$$
$$\vdots$$
\Large
$$\frac{\partial \vec{u}_i}{\partial t} = \left(\frac{\nu}{\Delta x^2} - \frac{c}{2 \Delta x}\right)u_{i+1} - \left(\frac{2}{\Delta x^2}\right) u_i + \left( \frac{\nu}{\Delta x^2} + \frac{c}{2 \Delta x} \right) u_{i-1}$$
\LARGE
\subsubsection{Viscous Burger's Equation}
Newton's Method:
$$\mathrm{\textbf{J}} \left( u^{n+1} - u^n \right) - \mathrm{\textbf{R}} \left( u^n \right)=0$$

\subsection{Reduced Order Model}

The goal of ROM is to lower the rank of the matrices involved in evolving the system; this can be done with \textbf{singular value decomposition}.
$$\mathrm{min} \{ ||\textbf{X} - \tilde{\textbf{X}}||_F \} \ni \mathrm{rank}(\tilde{\textbf{X}}) \ll \mathrm{rank}(\textbf{X})$$
$$\tilde{\textbf{X}} = U \Sigma V^T$$
The number of modes we use to create the model is dependent on how many singular values are used. In the case of linear convection and diffusion, we used five modes to model with nearly no error.

%----------------------------------------------------------------------------------------
%	RESULTS 
%----------------------------------------------------------------------------------------
\vspace{-2cm}
\section{Results}

\includegraphics[width=\linewidth]{fig.pdf}

Preliminary timing tests show a speed-up of about 50\%.

%----------------------------------------------------------------------------------------
%	CONCLUSIONS
%----------------------------------------------------------------------------------------

%\color{Navy} % SaddleBrown color for the conclusions to make them stand out

%\section{Conclusions}

%\begin{itemize}
%\item 
%\end{itemize}

\color{DarkSlateGray} % Set the color back to DarkSlateGray for the rest of the content

%----------------------------------------------------------------------------------------
%	FORTHCOMING RESEARCH
%----------------------------------------------------------------------------------------
\vspace{-1.5cm}
\section{Forthcoming Research}

\begin{itemize}
	\item Model Black-Scholes Equation
	\item Merton Model
	$$\frac{\partial u}{\partial \tau} = \frac{1}{2} \sigma^2_s s^2 \frac{\partial^2 u}{\partial s^2} + (r-\mu \xi) s \frac{\partial u}{\partial s} - (r+\mu)u + \mu \int_{0}^{\infty}u (sy,\tau)p(y)dy$$
	\item Discrete empirical interpolation method
\end{itemize}
%----------------------------------------------------------------------------------------
%	REFERENCES
%----------------------------------------------------------------------------------------
\large
\nocite{*} % Print all references regardless of whether they were cited in the poster or not
\bibliographystyle{plain} % Plain referencing style
\bibliography{poster} % Use the example bibliography file sample.bib

%----------------------------------------------------------------------------------------
%	ACKNOWLEDGEMENTS
%----------------------------------------------------------------------------------------

\section{Acknowledgements}

\begin{minipage}[c]{0.5\linewidth}
\begin{centering}
\includegraphics[width=\linewidth]{pure_logof.pdf}
\end{centering}
\end{minipage}
\begin{minipage}[c]{0.5\linewidth}
\begin{centering}
\includegraphics[width=\linewidth]{logo.jpg}
\end{centering}
\end{minipage}
%----------------------------------------------------------------------------------------

\end{multicols}
\end{document}